\chapter{Entropy Shaping as the Core Mechanism}

Delta--sigma modulation is traditionally viewed as a mechanism that shapes quantization noise by embedding a quantizer inside a feedback loop containing one or more integrators. In classical texts, noise shaping is described primarily in the frequency domain. The noise transfer function redistributes the spectral density of quantization error so that in-band components are suppressed and out-of-band components are amplified. While this spectral description is accurate, it conceals a deeper structural truth: noise shaping is entropy shaping. The delta--sigma loop manipulates the entropy field generated by quantization, directing uncertainty into spectral regions where it is harmless while preserving coherence in the region of interest.

The purpose of this chapter is to introduce the entropy field ( S(x,t) ), to show how it interacts with the loop filter and quantizer, and to demonstrate that delta--sigma modulation is fundamentally an entropy-transport process. This reinterpretation will provide the conceptual and mathematical foundations for the flow-dynamic, manifold-aligned, and derived-geometric analyses developed in the subsequent chapters of this part.

\section*{4.1 Quantization as Localized Entropy Injection}

Quantization converts a continuous value into a discrete symbol. This collapse operation, denoted ( Q : \mathbb{R} \to \Lambda ), is a form of coarse-graining. In information-theoretic terms, it compresses the infinite precision of the input into a finite representation, thereby discarding information. This loss of information corresponds to an increase in entropy. Quantization therefore injects entropy into the system at each time step.

Let ( u[n] ) be the input and ( y[n] = Q(v[n]) ) the quantized output. The quantization error
[
e[n] = u[n] - y[n]
]
contains the entropy injected at time (n). In classical analyses, ( e[n] ) is treated as white noise under suitable assumptions \citep{widrow1956statistical}. However, from a field-theoretic perspective, ( e[n] ) defines the local update of an entropy field. To formalize this, consider that the internal state of a delta--sigma modulator evolves according to
[
x[n+1] = f(x[n]) + L(e[n]),
]
where (L) is the loop filter. At each step, the quantizer injects an amount of entropy proportional to the unpredictability of (y[n]) given the state (v[n]). Let
[
S[n] = H(y[n] \mid v[n])
]
denote the conditional entropy of the quantizer output given the internal signal. Although the quantizer is deterministic, the effective entropy arises from the uncertainty of the internal signal relative to the partition induced by the quantization thresholds.

Thus, the quantization operation may be regarded as a localized entropy impulse ( \Delta S[n] ) applied to the system. Unlike perceived noise, this entropy is not an external disturbance but an intrinsic consequence of symbolic collapse.

\section*{4.2 The Entropy Field (S(x,t))}

To describe how this injected entropy evolves within the loop, we introduce the entropy field
[
S : \mathcal{X} \times \mathbb{Z} \to \mathbb{R}_{\ge 0},
]
assigning to each state value ( x[n] \in \mathcal{X} ) an entropy density ( S(x[n],n) ). This field is discrete in time but may be continuous in state. The function ( S(x,t) ) plays the role of a scalar field similar to entropy densities in nonequilibrium thermodynamics \citep{callen1985thermodynamics} and to the scalar entropy fields that appear in the RSVP framework.

Using this field, the quantizer injects entropy according to
[
S(x[n],n) = S(x[n],n^-)+\Delta S[n],
]
where (n^-) denotes the instant before quantization. The loop filter then redistributes this entropy spatially (across the state manifold) and spectrally (across frequency modes), according to its vector-field structure.

\section*{4.3 Oversampling as Entropic Dilution}

Oversampling reduces the local entropy gradient. Classical analyses describe oversampling as sampling faster than the Nyquist rate to spread quantization noise across a wider frequency range \citep{candy1992oversampling}. From the entropy-field perspective, the effect is more fundamental. Oversampling increases the temporal resolution at which the entropy field is sampled. Given a fixed entropy injection per sample, increasing the number of samples per unit time decreases the entropy per unit frequency band.

Let the oversampling ratio be ( K ). Consider the entropy per Nyquist interval:
[
S_{\text{Nyq}} = \frac{1}{K} \sum_{i=0}^{K-1} \Delta S[n+i].
]
For fixed quantizer precision, the average entropy density decreases linearly with oversampling. Thus, oversampling does not merely dilute noise; it smooths the entropy field, decreasing the amplitude of entropy gradients that the loop filter must correct. This interpretation yields a more direct understanding of why oversampling improves stability margins.

\section*{4.4 Loop Filtering as Entropy Transport}

The loop filter behaves as a vector field acting on the entropy field. Consider the classical form
[
x[n+1] = x[n] + e[n]
]
for a first-order modulator. In the absence of quantization, (x[n]) would integrate the input. When quantization injects entropy, the integrator transports this entropy to the internal state. Because the quantization error is correlated across time, the transported entropy accumulates preferentially in certain spectral modes.

Generalizing to higher-order loop filters, let the loop operator be represented in the continuous-time limit as a differential operator ( \mathcal{L} ):
[
\frac{dx}{dt} = \mathcal{L}(e(t)).
]
Then the entropy field evolves according to
[
\frac{\partial S}{\partial t} = -\nabla \cdot (\mathbf{v} S) + D \Delta S + \Delta S_{\text{inj}},
]
where ( \mathbf{v} ) is the vector field generated by the loop filter and (D) is a diffusion coefficient representing smoothing. This is a conservation equation with injection and advection. In delta--sigma language, ( \mathbf{v} ) is determined by the loop filter transfer function, and ( D ) corresponds to the implicit smoothing provided by higher-order integration.

Thus, classical noise shaping is a special case of an entropy-transport PDE. The noise transfer function specifies the spectral redistribution, while the loop dynamics specify the vector flow.

\section*{4.5 Boundary Conditions and the Role of the Quantizer}

Entropy transport in a delta--sigma loop is constrained by the nonlinear boundary conditions imposed by the quantizer. Each quantization cell corresponds to a region of the state space where the entropy injection is locally constant. The boundaries between cells are surfaces of discontinuity. When the internal state crosses a boundary, the quantizer output changes abruptly, producing a sharp entropy impulse.

Let the quantization thresholds be ( {\tau_k} ). The entropy injection rate satisfies
[
\Delta S[n] = S_{k^+} - S_{k^-}
]
when (x[n]) crosses from cell (k^-) to cell (k^+). These discontinuities act as entropy sources that modulate the divergence of the vector field:
[
\nabla \cdot \mathbf{v}(x,t) = -\frac{dS}{dt} + R(x,t),
]
where ( R(x,t) ) captures geometric and nonlinear corrections. In classical analyses, these contributions are effectively ignored by linearizing the quantizer as an additive noise source. In the entropy-field interpretation, they become essential to understanding stability, limit cycles, and chaotic regimes.

\section*{4.6 Entropy Shaping: The Core Mechanism}

We can now state the central thesis of this part of the book:

**Delta--sigma modulation is an entropy-shaping dynamical system. Quantization injects localized entropy into the state. The loop filter transports and redistributes this entropy, concentrating it in high-frequency modes where it is harmless while preserving low-entropy coherence in the band of interest.**

The classical noise transfer function captures only the frequency-domain manifestation of this process. The geometric interpretation captures its full dynamical structure. Entropy shaping is not a metaphor but a mathematically precise description of the internal processes of the delta--sigma loop. In the remainder of Part~II, this viewpoint will be developed using manifold geometry, vector-field dynamics, and derived structures inspired by the RSVP framework.

\section*{4.7 Transition to the Flow-Dynamic Interpretation}

This chapter has introduced the entropy field, the vector-field behavior of loop filters, and the PDE-like structure underlying delta--sigma feedback. In the next chapter, we develop the flow-dynamic interpretation in full detail, showing how the internal state variables of the modulator behave as vector-field components analogous to baryon-like flows in RSVP theory. The notions of divergence, stability, and attractors will be recast in geometric terms. This will complete the foundation for manifold-aligned representations (Chapter 6) and the deeper reformulation of delta--sigma modulation as a derived geometric object (Chapter 9 and beyond).
